\SGAProof{Proof of \SGARefLink{sga3:xvii:theorem:6:1:1:item:C}{6.1.1(C)}}
If \(U_1'\) is a lift in \(E\) of an algebraic subgroup \(U_1\)
of \(U\), then \(U_1'\) is invariant in \(E\), since \(E\) is
assumed commutative, and we may pass to the quotient by \(U_1'\).
Using a composition series of \(U\) (\SGARefLink{sga3:xvii:theorem:3:5:item:ii}{3.5(ii)}), the preceding
remark reduces us to the case in which \(U\) is smooth or equal
to \(\boldsymbol\alpha_p\).  But then \(E\) is trivial by
(A)(iv) and (B).

\SGASubsection{6.4}{Examples of nontrivial extensions of a
unipotent group \(U\) by a group \(H\) of multiplicative type}
\SGARefTarget{sga3:xvii:subsection:6:4}

In view of \SGARefLink{sga3:xvii:theorem:6:1:1:item:A}{6.1.1(A)}(iv), the question arises only in
characteristic \(p>0\).

\begin{enumerate}
\item[(a)] \SGARefTarget{sga3:xvii:subsection:6:4:item:a}
  \(H=\mathbb G_m\), and \(U\) is a nontrivial form of
  \(\mathbb G_a\) (cf.\ Appendix III, \S5).

  Let \(k\) be a nonperfect field of characteristic \(2\), and
  let \(u\in k\) be such that \(u^{1/2}\notin k\).  Consider the
  affine group \(E\) with coordinate ring
  \[
    k\bigl[X,Y,(X^2+uY^2)^{-1}\bigr],
  \]
  whose multiplication is given by the comorphism
  \[
    (X,Y)\longmapsto
      \bigl(XX'+uYY',\,XY'+YX'\bigr).
  \]
  The group \(E\) is smooth, connected, commutative, and of
  dimension \(2\); the subscheme \(Y=0\) defines a subgroup
  \(H\simeq\mathbb G_m\).  The kernel \(K\) of the squaring map
  on \(E\) has equation
  \[
    X^2+uY^2=1,
  \]
  and hence has dimension \(1\).  The group \(K\) contains the
  unipotent radical of \(E\), defined over \(\bar k\), as well as
  the contribution of \(H\), which is isomorphic to
  \(\boldsymbol\mu_2\).  Since \(K\) is reduced over \(\bar k\),
  the unipotent radical of \(E\) is not defined over \(k\), and
  \(U=E/H\) does not lift in \(E\).

  One checks immediately that \(U\) is the form of
  \(\mathbb G_a\) with coordinate ring
  \[
    k[V,W]/(V+uV^2+W^2),
  \]
  the morphism \(E\longrightarrow U\) corresponding to the
  comorphism
  \[
    V\longmapsto\frac{Y^2}{X^2+uY^2},
    \qquad
    W\longmapsto\frac{XY}{X^2+uY^2}.
  \]

\item[(b)] \SGARefTarget{sga3:xvii:subsection:6:4:item:b}
  \(H=\mathbb G_m\), \(U=\mathbb Z/2\mathbb Z\), and \(k\) is
  nonperfect of characteristic \(2\).

  Choose \(k\) and \(u\) as in \SGARefLink{sga3:xvii:appendix:A}{(a)}, and consider the subgroup
  \(E\) of \(\mathrm{GL}_2\) generated by the element
  \[
    X=
    \begin{pmatrix}0&1\\u&0\end{pmatrix}
    =
    \begin{pmatrix}u^{1/2}&0\\0&u^{1/2}\end{pmatrix}
    \begin{pmatrix}0&u^{-1/2}\\u^{1/2}&0\end{pmatrix}.
  \]
  The group \(E\) is an extension of \(\mathbb Z/2\mathbb Z\)
  by \(\mathbb G_m\), but this extension is not trivial because
  the unipotent part of \(X\) is not defined over \(k\).

\item[(c)] \SGARefTarget{sga3:xvii:subsection:6:4:item:c}
  \(H=\boldsymbol\mu_p\), \(U=\boldsymbol\alpha_p\), and \(k\)
  is nonperfect.

  Let \(\mathfrak e\) be the commutative restricted Lie algebra
  generated by two elements \(X\) and \(Y\) such that
  \[
    X^{(p)}=X,\qquad Y^{(p)}=aX.
  \]
  By Appendix II, \SGARefLink{sga3:ii:subsection:2-2}{2.2}, \(\mathfrak e\) is the restricted Lie
  algebra of an algebraic group \(E\) which is an extension of
  \(\boldsymbol\alpha_p\) by \(\boldsymbol\mu_p\).  This
  extension is trivial if and only if there exists \(b\in k\)
  such that \(b^p=a\), for then
  \[
    (Y-bX)^{(p)}=0.
  \]

\item[(d)] \SGARefTarget{sga3:xvii:subsection:6:4:item:d}
  \(H=\boldsymbol\mu_2\), \(U=\boldsymbol\alpha_2\times
  \boldsymbol\alpha_2\), \(E\) noncommutative, and \(k\) a field
  of characteristic \(2\).

  Consider the special linear group \(\mathrm{SL}_{2,k}\), and
  put \(E=F(\mathrm{SL}_{2,k})\).  The group \(E\) is a radicial
  group of height \(1\), whose Lie algebra is generated by three
  elements \(X,Y,Z\) satisfying
  \[
    [X,Y]=Z,\qquad [X,Z]=[Y,Z]=0,
  \]
  \[
    X^{(p)}=Y^{(p)}=0,\qquad Z^{(p)}=Z.
  \]
\end{enumerate}
